\section{Exact counting and Poisson reduction}
\label{sec:exact}
\subsection{A formal generating function and a graph model}
Expanding the finite binomial at each $m$ gives the formal identity
\begin{equation}\label{eq:ogf}
 \sum_{n\ge0}a_p(n)x^n
       =\sum_{m\ge0}x^m\bigl(1+(mx)^p\bigr)^m.
\end{equation}
Indeed, the term indexed by $k$ on the right is
$\binom mk m^{pk}x^{m+pk}$. Setting $m=n-pk$ gives
\eqref{eq:counts}, and $k\le m$ becomes $qk\le n$.
Every coefficient receives only finitely many contributions, so no analytic
convergence is being claimed. In fact, the coefficients have superexponential
growth.

For a direct interpretation, use the canonical labeled vertex set
$[m]=\{1,\ldots,m\}$. Choose $k$ active vertices. Each active vertex has $p$
ordered, or color-labeled, outgoing arc slots, and each slot independently
chooses a target in $[m]$. Loops and repeated targets are allowed.
Inactive vertices have no outgoing slots. There are
$\binom mk m^{pk}$ such structures. The size statistic is
\[
 n=m+pk=\hbox{number of vertices plus number of arcs}.
\]
The disjoint union over the possible $m$ therefore has size $a_p(n)$.
This is a labeled model with a nonstandard size statistic;
\eqref{eq:ogf} is not an ordinary unlabeled-digraph generating function.

The number of inactive vertices is
\begin{equation}\label{eq:rdef}
 r=m-k=n-qk,\qquad
 k=\frac{n-r}{q},\qquad m=\frac{n+pr}{q}.
\end{equation}
For a uniformly chosen structure of size $n$, denote this random number by
$R_n$. Its support is exactly
$\{r:0\le r\le n,\ r\equiv n\pmod q\}$.

\subsection{An exact change of variables}
For $n>0$ write $N=n/q$ and $B=N^{pn/q}$.
For an allowed $r$, use $\binom mk=\binom mr$ to get
\[
 m^{pk}\binom mk
  =\frac{m^{pk+r}}{r!}\prod_{h=0}^{r-1}\left(1-\frac hm\right).
\]
Since $pk+r=(pn+r)/q$ and $m=N(1+pr/n)$, this becomes
\begin{equation}\label{eq:exactterm}
 m^{pk}\binom mk=B\frac{\lambda^r}{r!}e^{\Delta_n(r)},
\end{equation}
where
\begin{equation}\label{eq:Delta}
 \Delta_n(r)=\frac{pn+r}{q}\log(1+pr/n)-\frac{p^2r}{q}
      +\sum_{h=0}^{r-1}\log\left(1-\frac{qh}{n+pr}\right).
\end{equation}
All factors inside the logarithms are positive for integers $0\le r\le n$:
$n+pr-q(r-1)=n-r+q>0$.
The empty product and empty sum at $r=0$ are one and zero respectively.
The same expression defines a real function at every integer in this range,
not only those satisfying the congruence.

If $X\sim\Pois(\lambda)$, summation gives the exact identity
\begin{equation}\label{eq:exactpoisson}
 q e^{-\lambda}\frac{a_p(n)}{B}
  =q\E\left[e^{\Delta_n(X)}\1_{\{X\le n\}}
                    \1_{\{X\equiv n\ (q)\}}\right].
\end{equation}
For $p\ge2$ the left side is $\Rcal_p(n)$.
For $p=1$ it is $e^{-d}a_1(n)/M_1(n)$.
The factor $q$ is the lattice spacing of the residue condition. It does
not come from a Stirling approximation; no approximation has yet been used.

\subsection{Poisson polynomial identities}
Let $T_j$ denote the Touchard polynomial,
\begin{equation}\label{eq:touchard}
 T_j(z)=\sum_{\ell=0}^j\left\{\begin{matrix}j\\\ell\end{matrix}\right\}z^\ell,
 \qquad \E X^j=T_j(\lambda).
\end{equation}
Here the braces are Stirling numbers of the second kind.
This follows by writing $X^j$ in the falling-factorial basis and using
$\E(X)_\ell=\lambda^\ell$.
In particular, for each fixed $j$,
$\E X^j=O_j((1+\lambda)^j)$.

Put $\omega=e^{2\pi\ii/q}$ and
$\kappa_q=1-\cos(2\pi/q)>0$.
For any integer residue $a$ and any nonnegative integer $j$,
\begin{equation}\label{eq:filter}
 q\E[X^j\1_{\{X\equiv a\ (q)\}}]
  =\sum_{\ell=0}^{q-1}\omega^{-a\ell}
       e^{\lambda(\omega^\ell-1)}T_j(\lambda\omega^\ell).
\end{equation}
To see this, insert
$q\1_{\{X\equiv a\}}=\sum_{\ell=0}^{q-1}\omega^{\ell(X-a)}$
and use the Poisson generating function. The $\ell=0$ term is $T_j(\lambda)$;
all other terms are
$O_{q,j}((1+\lambda)^j e^{-\kappa_q\lambda})$.
For a fixed polynomial, residue restriction therefore changes its Poisson
average only by an exponentially small absolute error.
This statement concerns polynomial averages. It does not identify a
relative secondary exponential sector of the full nonpolynomial sum.

\subsection{Exact monotonicity}
\begin{proposition}\label{prop:monotone}
For every $p\ge1$,
\[
 a_p(0)=\cdots=a_p(p)=1,\qquad a_p(p+1)=2,
\]
and $a_p(n+1)>a_p(n)$ for every $n\ge p$.
Consequently, for each real $y>1$ the threshold
\begin{equation}\label{eq:threshold}
 T_p(y)=\min\{n\ge0:a_p(n)\ge y\}
\end{equation}
is well-defined.
\end{proposition}
\begin{proof}
For each previously allowed $k$, increasing $n$ increases
$m=n-pk\ge k$. Both $m^{pk}$ and $\binom mk$ are nondecreasing in $m$,
and all newly allowed summands are positive. At the first step $n=p$,
the new $k=1$ term gives strictness. Thereafter the common $k=1$ term
$(n-p)^{p+1}$ strictly increases. It also proves unboundedness.
The initial equalities are immediate from \eqref{eq:counts}.
\end{proof}
